% TOP_PROBLEM: {"id": 8400041, "title": "O36 — Integer multiplication in linear bit complexity", "category_id": 3, "category": {"id": 3, "name": "graph_theory", "display_name": "Graph Theory", "description": "Problems involving graphs, networks, and their properties.", "slug": "graph-theory", "order_index": 3, "created_at": "2026-07-31T15:26:25.671Z"}, "status": "partially_solved", "problem_number": "AMR-083-0041", "set_id": 13}
% FIRST_POSTED: 2026-10-01
% ENTRY_KIND: statement_only
% UnsolvedMath ID 8400041; source catalogue code AMR-083-0041
% !TeX program = lualatex
% Definition and sources only; this file is not an LLM solution attempt.
\documentclass[11pt,a4paper]{article}
\usepackage[margin=25mm]{geometry}
\usepackage{fontspec}
\setmainfont{lmroman10-regular.otf}[BoldFont=lmroman10-bold.otf,
 ItalicFont=lmroman10-italic.otf,BoldItalicFont=lmroman10-bolditalic.otf]
\usepackage{amsmath,amssymb,mathtools}
\usepackage{unicode-math}
\setmathfont{latinmodern-math.otf}
\AtBeginDocument{\let\setminus\smallsetminus}
\usepackage{xurl,hyperref}
\hypersetup{colorlinks=true,urlcolor=blue}
\setcounter{secnumdepth}{0}
\setlength{\parindent}{0pt}
\setlength{\parskip}{5pt}
\setlength{\emergencystretch}{3em}
\begin{document}
\section*{O36 — Integer multiplication in linear bit complexity}\label{8400041}
{\small UnsolvedMath ID 8400041 · AMR-083-0041 · Graph Theory · AMR Open Problem Lists}
\subsection{Definitions and mathematical statement}
Let $a,b$ be positive integers represented by their ordinary binary expansions, with lengths $m$ and $n$. Set $L=m+n$. Their product is to be output exactly in binary, including every digit; its length is at most $L$. Since $L$ differs from $\log_2(ab)$ by a bounded additive amount, linear time in $L$ is the intended meaning of the source's $O(\log(ab))$ bound.

Does a deterministic classical algorithm multiply every such pair in at most $CL$ elementary bit operations for some absolute constant $C$? Use the standard multitape Turing-machine bit model, or a comparably explicit bit-cost model in which processing arbitrarily long integer words is not a single free operation. The algorithm must be uniform over all input lengths.

The requirement covers unbalanced lengths as well as equal-length inputs; padding reduces these conventions to equivalent linear bounds. It asks for the exact full product, not a residue, a numerical approximation, or only its leading digits. Reading the inputs and writing the answer count as computation.

An algorithm whose cost includes an unbounded factor such as $\log L$ does not by itself establish the stated linear bound, however slowly that factor grows. Conversely, an impossibility result would need to address the stated bit model rather than a more restricted multiplication procedure. The computational model is essential to making the linear-time target unambiguous.
\subsection{Short English statement}
Can exact multiplication of binary integers be performed in time linear in their total number of bits?
\subsection{Sources}
\begin{itemize}
\item[S1] ULAM.AI and the UnsolvedMath Contributors, supplied problems.json, record 8400041 (AMR-083-0041), AMR Open Problem Lists. Catalogue identity and source transcription; CC BY 4.0.
\url{https://huggingface.co/datasets/ulamai/UnsolvedMath}.
\item[S2] Adleman \& McCurley - Open Problems in Number Theory Complexity II (1994), O36, p23 / LNCS p313. Original source indexed by AMR-083-0041.
\url{https://mccurley.org/papers/open.ps.gz}.
\end{itemize}
\end{document}
